\documentclass[11pt,letterpaper]{article}

% ---- Packages ----
\usepackage[margin=1in]{geometry}
\usepackage{amsmath,amssymb,amsthm}
\usepackage{bm}
\usepackage{booktabs}
\usepackage{xcolor}
\usepackage[colorlinks=true,linkcolor=blue,citecolor=blue,urlcolor=blue]{hyperref}
\usepackage{graphicx}

% ---- Math commands ----
\newcommand{\ind}[1]{\mathbf{1}\{#1\}}
\newcommand{\E}{\mathbb{E}}
\newcommand{\V}{\mathbb{V}}
\newcommand{\R}{\mathbb{R}}
\newcommand{\trt}{\text{t}}
\newcommand{\ctrl}{\text{c}}
\newcommand{\pr}{\text{pr}}
\newcommand{\tp}{\text{tx}}
\newcommand{\cp}{\text{cx}}
\newcommand{\Prm}[1]{\widetilde{#1}}
\newcommand{\sgn}{\operatorname{sgn}}

\title{\vspace{-1.5cm}%
  \textbf{Stratified RMST Boosting for Subgroup Identification}\\[4pt]
  \large Augmenting the Value Function with Prognostic Strata}
\author{Yue Shentu \& Natasha}
\date{May 28, 2026\\[6pt]\small Revised v2: corrected gradient scaling}

\begin{document}
\maketitle
\tableofcontents
\newpage

% ======================================================================
\section{Introduction}
% ======================================================================

This document defines a stratified version of RMST-based subgroup boosting
(``RMST CAVBoost'').  The original method learns a soft subgroup
assignment $p_i \in [0,1]$ (the probability that patient $i$ belongs to the
treatment-preferring subgroup) by optimizing a value function through
gradient boosting trees.  The proposed refinement replaces the original
pooled value function with one computed within prognostic strata and
aggregated.

We focus on \textbf{stratified training}, which replaces the value
function used during boosting with the stratified version.  The
gradient is modified to incorporate stratum-specific RMST differences.
This is the version evaluated in the simulations below.

% ======================================================================
\section{Notation}
% ======================================================================

For each subject $i = 1,\dots,N$, we observe:
\begin{itemize}
  \item $\bm{Z}_i$: $p$-vector of baseline covariates;
  \item $A_i \in \{0,1\}$: treatment (1 = active, 0 = control);
  \item $U_i = \min(T_i, C_i, \tau)$: observed time;
  \item $\widetilde\delta_i = \ind{\min(T_i,\tau) \le C_i}$: event within
        $\tau$ before censoring.
\end{itemize}

The restricted mean survival time (RMST) at horizon $\tau$ is
$\mu(\tau) = \int_0^\tau S(t)\,dt$, estimated by the area under the
Kaplan--Meier curve.

% ======================================================================
\section{Prognostic Stratification}
% ======================================================================

A prognostic score $s_i$ is a univariate summary of $\bm{Z}_i$ that
predicts survival under the control condition.  Because treatment is
randomized, we can use all patients (both arms) to estimate $s_i$ as long
as treatment is excluded from the predictor set.

We use a cross-fitted elastic net Cox model to estimate $s_i$,
balancing the need to handle $p \approx n/10$ covariates while
avoiding overfit within each fold.  The implementation:
\begin{enumerate}
  \item Partition the training data into $F = 5$ random folds.
  \item For fold $f = 1,\dots,F$:
    \begin{enumerate}
      \item Fit $\texttt{cv.glmnet}(\bm{Z}, \text{Surv}(T,C),$
            $\ \texttt{family="cox"}, \texttt{alpha=0.5})$ on the
            out-of-fold observations.
      \item Predict the linear predictor $s_i$ for in-fold
            observations at $\texttt{lambda.min}$.
    \end{enumerate}
  \item Collect the out-of-fold predictions $\widehat{s}_i$ for all
        training patients.
  \item Define strata by $K$ quantiles of $\widehat{s}_i$.
\end{enumerate}
Cross-fitting ensures that the stratum assignment for each patient is
independent of its own data, eliminating the training-set overlap bias
that would arise from fitting and predicting on the same observations.
The elastic net penalty (mixing parameter $\alpha=0.5$, cross-validated
$\lambda$) automatically selects among the 52 covariates and shrinks
noisy coefficients, which is essential when the Cox model is
proportional-hazards misspecified (e.g., under the quadratic prognostic
structure of Zhang scenarios~S4--S6).

Stratify $s_i$ into $K$ groups by its quantiles.  Let $\mathcal{S}_k$
be the index set of stratum $k$, with $n_k = |\mathcal{S}_k|$ patients.
For $K = 4$, these are quartiles; $K = 6$ and $K = 8$ are used in
sensitivity analyses.

% ======================================================================
\section{Stratified Training}
\label{sec:strat-train}
% ======================================================================

In this mode, the value function used during XGBoost training is itself
stratified.  This changes the gradient that each tree fits, potentially
altering the learned $p_i$.

\subsection{Value Function}

The stratified value function mirrors the original CAVBoost structure
exactly, but applied within each stratum.  Recall that the original
value function is:
\[
V_{\text{orig}} = \Bigl(\sum_i p_i\Bigr)\widehat\Delta^{(1)}
                - \Bigl(\sum_i (1-p_i)\Bigr)\widehat\Delta^{(2)},
\]
where $\widehat\Delta^{(1)}$ uses $p_i$-weighted KM (subgroup 1,
treatment-preferring) and $\widehat\Delta^{(2)}$ uses $(1-p_i)$-weighted
KM (subgroup 2, control-preferring).

The stratified version applies the same logic within each stratum $k$.
Define within-stratum subgroup~1 KM (p-weighted) and subgroup~2 KM
((1-p)-weighted), giving two RMST differences per stratum:

\begin{align*}
d_k^{(1)} &=
  \widehat\mu_{1,k}^{(1)}(\tau) - \widehat\mu_{0,k}^{(1)}(\tau),\\
d_k^{(2)} &=
  \widehat\mu_{1,k}^{(2)}(\tau) - \widehat\mu_{0,k}^{(2)}(\tau),
\end{align*}
where $\widehat\mu_{a,k}^{(1)}(\tau)$ is the $p$-weighted KM RMST within
stratum $k$, arm $a$, and $\widehat\mu_{a,k}^{(2)}(\tau)$ is the
$(1-p)$-weighted equivalent.

The stratified value function is:

\[
\boxed{
V_{\text{strat}}(\bm{p}) =
\sum_{k=1}^K W_k\, d_k^{(1)}
\;\;\;\;\;\;\;\;-
\;\;\;\;\;\;\;\;
\sum_{k=1}^K (n_k - W_k)\, d_k^{(2)}
}
\]
where $W_k = \sum_{i\in\mathcal{S}_k} p_i$ and $n_k = |\mathcal{S}_k|$.

This is the same \emph{difference of deltas} format as the original, but
the KM computations are restricted to each stratum.  No normalization by
$N$ is applied, matching the scale convention of the original.

\subsection{Analytic Gradient}

The gradient mirrors the original CAVBoost derivation.  For patient $j$
in stratum $k$, taking the derivative of each term in $V_{\text{strat}}$
with respect to $\eta_j = \log(p_j/(1-p_j))$:

\textbf{Term 1:} $\displaystyle\frac{\partial}{\partial\eta_j}\bigl[W_k\,d_k^{(1)}\bigr]
= p_j(1-p_j)\,d_k^{(1)} + W_k\,\frac{\partial d_k^{(1)}}{\partial\eta_j}$

\textbf{Term 2:} $\displaystyle-\frac{\partial}{\partial\eta_j}\bigl[(n_k-W_k)\,d_k^{(2)}\bigr]
= p_j(1-p_j)\,d_k^{(2)} - (n_k-W_k)\,\frac{\partial d_k^{(2)}}{\partial\eta_j}$

The chain rule $\partial d_k^{(g)}/\partial\eta_j =
p_j(1-p_j)\,\partial d_k^{(g)}/\partial p_j$ gives the full gradient:

\[
\boxed{
\frac{\partial V_{\text{strat}}}{\partial\eta_j}
= p_j(1-p_j)\Bigl[
  d_k^{(1)} + d_k^{(2)}
  + W_k\,\frac{\partial d_k^{(1)}}{\partial p_j}
  - (n_k-W_k)\,\frac{\partial d_k^{(2)}}{\partial p_j}
\Bigr].
}
\]

The XGBoost loss function minimizes, so the gradient passed to the tree
learner is $g_j = -\partial V_{\text{strat}}/\partial\eta_j$.

\paragraph{The influence functions.}
The terms $\partial d_k^{(1)}/\partial p_j$ and
$\partial d_k^{(2)}/\partial p_j$ are the derivatives of the
within-stratum $p$-weighted and $(1-p)$-weighted RMST differences with
respect to patient $j$'s subgroup probability.  They are computed by the
same influence-function machinery as the original CAVBoost (the
$\text{hazard\_inc}$ function in the R code), evaluated on the events
within stratum $k$ only.  The per-stratum computation is $K$ times
faster than the full-sample version because each stratum has roughly
$n_k \approx N/K$ patients.

\paragraph{Block-diagonal structure.}
A critical property: $\partial d_k^{(g)}/\partial p_j = 0$ for $j \notin
\mathcal{S}_k$, because patient~$j$ does not contribute to the KM within
stratum~$k$.  This means the gradient for patient~$j$ depends only on
patients in the same stratum.  The XGBoost regression tree must therefore
learn stratum-relevant splits from the feature vector $\bm{Z}_j$, relying
on the fact that the prognostic score $s_i$ (used to define strata) is a
function of $\bm{Z}_i$.

\paragraph{Connection to the original.}
The original CAVBoost gradient (from \texttt{rmst\_cavboost\_clean.R}) is:
\[
g_j^{\text{(orig)}} = -p_j(1-p_j)\Big[
\widehat\Delta^{(1)} + \Bigl(\sum_i p_i\Bigr)\frac{\partial
  \widehat\Delta^{(1)}}{\partial p_j}
+ \widehat\Delta^{(2)} - \Bigl(\sum_i (1-p_i)\Bigr)\frac{\partial
  \widehat\Delta^{(2)}}{\partial p_j}
\Big].
\]
The stratified version is identical in structure, but with each term
replaced by its stratum-level analogue:
\begin{align*}
\widehat\Delta^{(1)} &\;\leftrightarrow\; d_k^{(1)},
& \widehat\Delta^{(2)} &\;\leftrightarrow\; d_k^{(2)},\\
\sum_i p_i &\;\leftrightarrow\; W_k,
& \sum_i (1-p_i) &\;\leftrightarrow\; n_k-W_k,\\
\frac{\partial\widehat\Delta^{(1)}}{\partial p_j}
&\;\leftrightarrow\; \frac{\partial d_k^{(1)}}{\partial p_j},
&
\frac{\partial\widehat\Delta^{(2)}}{\partial p_j}
&\;\leftrightarrow\; \frac{\partial d_k^{(2)}}{\partial p_j}.
\end{align*}
The implementation is a direct modification of the original: the
KM/hazard-increment computation loops over strata instead of the full
sample.

\subsection{Implementation}

The algorithm mirrors the original CAVBoost, but the KM and influence
computations loop over strata:
\begin{enumerate}
  \item Pre-compute cross-fitted prognostic score $\widehat{s}_i$ via
        $5$-fold elastic net Cox (glmnet, $\alpha=0.5$) and stratum
        assignments $\mathcal{S}_k$ by $K$ quantiles.  These are fixed
        throughout boosting.
  \item Initialize $\eta_i = 0$ ($p_i = 0.5$).
  \item For each boosting round:
        \begin{enumerate}
          \item Compute $p_i = 1/(1+e^{-\eta_i})$.
          \item For each stratum $k$:
            \begin{enumerate}
              \item Compute $p$-weighted KM within $\mathcal{S}_k$,
                    get $d_k^{(1)}$ and influence $\partial d_k^{(1)}/\partial p_j$.
              \item Compute $(1-p)$-weighted KM within $\mathcal{S}_k$,
                    get $d_k^{(2)}$ and influence $\partial d_k^{(2)}/\partial p_j$.
              \item Assemble gradient for patients in $\mathcal{S}_k$:
                    $g_j = -p_j(1-p_j)[d_k^{(1)} + d_k^{(2)}
                    + W_k\,\partial d_k^{(1)}/\partial p_j
                    - (n_k-W_k)\,\partial d_k^{(2)}/\partial p_j]$.
            \end{enumerate}
          \item Fit a regression tree to $\{g_j, \bm{Z}_j\}$.
          \item Update $\eta_i \leftarrow \eta_i + \eta_{\text{boost}}
                \times \text{tree}(\bm{Z}_i)$.
        \end{enumerate}
  \item After training: $\widehat G_i = \ind{p_i > 0.5}$.
\end{enumerate}

This is a direct adaptation of the original CAVBoost gradient, replacing
the global weighted~KM with per-stratum weighted~KM.  The R functions
from \texttt{rmst\_cavboost\_clean.R}---in particular the hazard
increment and sorted-data machinery---are reused without modification;
only the loop structure changes.





% ======================================================================
\section{Simulation Results}
\label{sec:results}
% ======================================================================

All classification metrics are computed on a held-out test set of the
same size as the training set.  Three methods are compared:
original CAVBoost, stratified CAVBoost ($K=4$, cross-fitted elastic net
Cox prognostic score), and virtual twin (separate
survival random forests (ranger, 200 trees,
min node size~10, RMST difference via trapezoidal integration).

\subsection{Complex Boundary Scenarios (Zhang et al. 2020)}

\label{sec:complex}

The Zhang et al. (2020) SubgroupBoost paper defines four scenarios with
increasing boundary complexity, using 52 covariates (50 noise, 2
predictive $S_1, S_2$), $N_{\text{train}} = 500$, $N_{\text{test}} = 2000$,
$\tau = 30$, and Setting~2 prognostic strength
($\beta_{Z_{1:4}} = 0.4$).  The prognostic score is estimated via
5-fold cross-fitted Cox PH, so the within-stratum gradient computation
uses out-of-fold predictions and avoids training-set overlap.

Results from 30 Monte Carlo reps:

\begin{center}
\begin{tabular}{lcccc}
\toprule
 & \multicolumn{3}{c}{Hold-out AUC} & Gain \\
\cmidrule{2-4}
Scenario (boundary) & Original & \textbf{Stratified} & VT & (Orig $\to$ Strat) \\
\midrule
S1: Linear ($S_1 = 0$) & 0.757 & \textbf{0.880} & 0.956 & $+$0.123\\
S3: U-shaped ($S_1$) & 0.844 & \textbf{0.932} & 0.509 & $+$0.088\\
S4: Enclave ($S_1 \times S_2$) & \textbf{0.778} & 0.764 & 0.508 & $-$0.014\\
S5: S-shaped ($S_1$, disjoint) & 0.805 & \textbf{0.883} & 0.595 & $+$0.078\\
\bottomrule
\end{tabular}
\end{center}

The key findings are:
\begin{itemize}
  \item \textbf{Virtual twin collapses on nonlinear boundaries} (S3: 0.509,
        S4: 0.508, S5: 0.595).  The per-arm linear Cox PH cannot capture
        non-monotone or interaction-based treatment effect surfaces.
  \item \textbf{Stratified CAVBoost excels when the subgroup boundary is
        smooth relative to prognosis} (S1: +12.3 AUC points, S3: +8.8,
        S5: +7.8).  The cross-fitted prognostic stratification aligns with
        smoothly varying treatment effects while the out-of-fold prediction
        eliminates the training-set overlap bias.
  \item \textbf{Stratification is neutral on disjoint boundaries}
        (S4 Enclave: $-$0.014, within Monte Carlo error).  When the true
        subgroup consists of disconnected regions in covariate space, the
        Cox PH-based prognostic score cannot capture the structure, but
        the original global gradient remains competitive.
  \item \textbf{Original CAVBoost is robust} across all scenarios
        (0.76--0.84 AUC), making it a reliable all-purpose baseline.
\end{itemize}

These results suggest that cross-fitted stratified CAVBoost is
consistently beneficial on smooth or monotone boundaries and never
catastrophically worse than the original.  The approach is
recommended as the default when the treatment-effect boundary is\,
believed to be continuous in the prognostic space.

\subsection{Gradient Combination Attempts}

Two approaches to combining the original and stratified gradients
were evaluated: alternating rounds (original on odd rounds, stratified
on even) and variance-gated blending ($\alpha_k$ derived from
Greenwood variance per stratum).  Both failed on the complex boundary
scenarios (AUC~0.55--0.60), because the two gradient functions are
minimizing different estimands and their per-observation directions
anti-correlate.  A holdout model-selection rule (train both models,
pick by held-out RMST gain) was also tested; it over-selected the
stratified model on S4 and S5 (where the original has a slight edge),
demonstrating that 150-patient holdouts lack the power to reliably
discriminate between the two methods.

\end{document}
